\chapter{Сквозной ток и условно индуцированные коэффициенты гравитации}
\label{ch:v10-throughflow-induced-gravity-coefficients}

\section{Хопфовский образ как стационарный канал}

Ранняя интуиция проекта представляла материю и вакуум как поток,
проходящий через хопфовский канал и возвращающийся наружу. Её строгий остаток
не требует буквальной трубки. Пусть
\begin{equation}
 \dot N=J_{\rm in}-J_{\rm out},
 \qquad \Theta=J_{\rm in}+J_{\rm out}.
 \label{eq:v10-throughflow-balance}
\end{equation}
При $J_{\rm in}=J_{\rm out}=J>0$ число накопленных ячеек постоянно, но
пропускная активность равна $\Theta=2J$. Таким образом, отсутствие роста
запаса не означает отсутствия движения через носитель.

Для физического сквозного тока необходима неравновесная сила $F>0$.
Производство энтропии равно
\begin{equation}
 \sigma=FJ>0.
 \label{eq:v10-throughflow-entropy-production}
\end{equation}
Именно $\sigma$, а не формальная сумма взаимно обратных мод, является
кандидатом на поддерживающий геометрический источник.

\section{Поддерживаемая и схлопнутая геометрия}

Введём вещественный параметр геометрического порядка $x$ и положительный
родитель
\begin{equation}
 \mathcal P_{\rm flow}(x)
 =\frac{\lambda}{4}\left(x^2-\frac{g\sigma}{\lambda}\right)^2,
 \qquad g,\lambda>0.
 \label{eq:v10-throughflow-order-parent}
\end{equation}
При включённом токе его устойчивые ветви удовлетворяют
\begin{equation}
 x_*^2=\frac{g\sigma}{\lambda},
 \qquad
 \mathcal P_{\rm flow}''(x_*)=2g\sigma>0,
 \label{eq:v10-throughflow-maintained-geometry}
\end{equation}
тогда как в начале координат
\begin{equation}
 \mathcal P_{\rm flow}''(0)=-g\sigma<0.
\end{equation}
Если ток выключен, то
\begin{equation}
 \mathcal P_{\rm flow}(x)\big|_{J=0}=\frac{\lambda}{4}x^4,
 \label{eq:v10-throughflow-collapse}
\end{equation}
и нарушенная ветвь исчезает. Это точное содержание образа: геометрия
поддерживается прохождением, а не накоплением вещества внутри канала.

\section{Условно индуцированные коэффициенты}

Положим $m=x_*^2=g\sigma/\lambda$. Тогда петлевой геометрический отклик
может иметь масштабно правильный вид
\begin{equation}
 A=\alpha m^2,
 \qquad B=\beta m.
 \label{eq:v10-throughflow-induced-coefficients}
\end{equation}
Однако кривизненный родитель по-прежнему выбирает только
\begin{equation}
 qm=\frac{\beta}{2\alpha}.
 \label{eq:v10-throughflow-curvature-invariant}
\end{equation}
Карта отношений между $(q,m,\sigma)$ имеет ранг $2$, ядро размерности $1$
с образующим $(1,-1,-1)^T$. Следовательно, совместное масштабирование
геометрии, тока и индуцированного семени остаётся незаметным.

\begin{theorem}[Сквозной ток условно поддерживает геометрию]
Равные ненулевые входящий и выходящий токи допускают постоянное число ячеек
при ненулевой пропускной активности. Положительное производство энтропии
создаёт устойчивую ненулевую ветвь $x_*^2=g\sigma/\lambda$, которая исчезает
при выключении тока.
\end{theorem}

\begin{nogocriterion}[Поддержание не равно абсолютной калибровке]
Сквозной ток объясняет, почему ненулевая геометрия может существовать только
динамически, но не создаёт секунду или длину, пока физически не выведены
величины $F$, $J$ и отношение $g/\lambda$. Их совместное масштабирование
оставляет точную нулевую моду.
\end{nogocriterion}

\begin{protocol}[Статус сквозного вакуумного тока]\begin{enumerate}
\item стационарный вход и выход: \textbf{построены};
\item ненулевая пропускная активность: \textbf{получена};
\item положительное производство энтропии: \textbf{условно получено};
\item поддерживаемая нарушенная геометрия: \textbf{получена};
\item архитектура: \textbf{$9/9$};
\item условное происхождение: \textbf{$4/4$};
\item физическое происхождение силы и тока: \textbf{$0/1$};
\item абсолютный гравитационный масштаб: \textbf{$0/1$};
\item следующий гейт: \textbf{аудит происхождения силы и сопротивления сквозного канала}.
\end{enumerate}\end{protocol}

Машинный сертификат сохранён в
\path{s2t_v10_cell_birth_four_volume_induced_gravity_coefficient_parent_origin_gate_results.json}.